\SourceLeaf{446}%
\begin{quote}
»Wenn von einem Punkte der krummen Fläche aus lauter kürzeste Linien von gleicher Länge gezogen sind, so machen diese mit der ihre Endpunkte verbindenden Linie überall rechte Winkel«.
\end{quote}
\noindent\hspace*{1.2em}Man kann auf ganz ähnliche Art beweisen, dass wenn auf der krummen Fläche irgend eine krumme Linie gegeben ist, und man von jedem Punkte derselben nach Einer Seite derselben und unter rechten Winkeln zu derselben lauter kürzeste Linien von gleicher Länge ausgehen lässt, deren Endpunkte durch eine Linie verbunden werden, diese von jenen überall unter rechten Winkeln geschnitten wird. Man braucht nur in der vorigen Entwickelung $\theta$ die Länge der gegebenen krummen Linie von einem beliebigen Punkte an bedeuten zu lassen, wo dann die vorigen Rechnungen ihre Gültigkeit behalten, nur dass die Richtigkeit des Werthes $u=0$ für $s=0$ jetzt schon in der Voraussetzung liegt.

\begin{center}18.\end{center}
\noindent\hspace*{1.2em}Die bei diesen Constructionen vorkommenden Grössenverhältnisse verdienen noch ausführlicher entwickelt zu werden. Wir haben also zuvörderst, wenn man Kürze halber $m$ für $\frac{\partial t}{\partial\theta}$ schreibt:
\begin{equation}
\frac{\partial x}{\partial s}=\xi,\qquad\frac{\partial y}{\partial s}=\eta,\qquad\frac{\partial z}{\partial s}=\zeta,\tag{1}
\end{equation}
\begin{equation}
\frac{\partial x}{\partial\theta}=m\xi',\qquad\frac{\partial y}{\partial\theta}=m\eta',\qquad\frac{\partial z}{\partial\theta}=m\zeta',\tag{2}
\end{equation}
\begin{equation}\xi\xi+\eta\eta+\zeta\zeta=1\tag{3}\end{equation}
\begin{equation}\xi'\xi'+\eta'\eta'+\zeta'\zeta'=1\tag{4}\end{equation}
\begin{equation}\xi\xi'+\eta\eta'+\zeta\zeta'=0.\tag{5}\end{equation}
Ferner
\begin{equation}XX+YY+ZZ=1\tag{6}\end{equation}
\begin{equation}X\xi+Y\eta+Z\zeta=0\tag{7}\end{equation}
\begin{equation}X\xi'+Y\eta'+Z\zeta'=0\tag{8}\end{equation}
und
\begin{equation}
\left\{\begin{aligned}X&=\zeta\eta'-\eta\zeta'\\Y&=\xi\zeta'-\zeta\xi'\\Z&=\eta\xi'-\xi\eta'\end{aligned}\right.\tag*{[9]}
\end{equation}
